\section{4. Experiments}

\subsubsection{Datasets}

To comprehensively evaluate the proposed Unified Dual-Mask Physical Model, we select four light field (LF) datasets that encompass a diverse range of surface reflectance properties, geometric complexities, and illumination conditions. The primary motivation for this selection is to cover the full spectrum of Lambertian, non-Lambertian (specular and scattering), and mixed urban scenes. This diversity is essential for validating the robustness of our model, particularly in challenging non-Lambertian regions where the photo-consistency assumption of traditional Epipolar Plane Image (EPI) based methods substantially breaks down. Note that the original HCI-Old dataset was excluded from our corpus due to format incompatibilities with our unified angular pipeline. The final dataset comprises 209 training scenes and 36 validation scenes.

\paragraph{Dataset Descriptions.}

\textbf{HCInew Dataset.} The HCInew dataset [1] is a widely recognized synthetic benchmark in LF depth estimation. It features complex indoor and outdoor scenes with rich textures and intricate geometric structures. The scenes predominantly exhibit Lambertian reflectance, providing a solid baseline for evaluating standard EPI slope extraction. The ground truth (GT) disparity maps are generated via ray-tracing engines with sub-pixel accuracy. 

\textbf{Wanner HCI Dataset.} Often referred to as the Wanner subset of the HCI benchmark [2], this dataset consists of synthetic scenes specifically designed to evaluate depth estimation algorithms under controlled lighting and geometric conditions. It primarily contains Lambertian surfaces with varying degrees of occlusion and textureless regions. The GT depth is derived directly from the synthetic rendering pipeline, ensuring perfect alignment with the sub-aperture images.

\textbf{Non-Lambertian Dataset (Zhenglong).} To explicitly address the limitations of existing methods on reflective surfaces, we incorporate a specialized non-Lambertian dataset curated by Zhenglong et al. [3]. This dataset synthesizes scenes incorporating complex Bidirectional Reflectance Distribution Functions (BRDFs), including high-gloss metals, dielectrics, and scattering materials (e.g., derived from the MERL BRDF database). The presence of specular highlights and view-dependent reflections severely violates the Lambertian assumption, causing significant EPI line distortions. The GT disparity maps represent the underlying geometric depth, strictly decoupled from the reflective layer.

\textbf{UrbanLF-Syn Dataset.} The UrbanLF-Syn dataset [4] provides large-scale synthetic urban environments, representing "Mixed" scenes that contain a heterogeneous combination of Lambertian (e.g., concrete, asphalt) and non-Lambertian (e.g., glass windows, metallic vehicle bodies) materials. This dataset is instrumental in evaluating the model's ability to seamlessly transition between different physical reflectance models within a single field of view. The GT depth maps are extracted from the high-fidelity 3D city models used for rendering.

\paragraph{Preprocessing and Sampling Strategy.}

To ensure a unified input representation and fair comparison across all datasets, we apply a standardized preprocessing pipeline:

1. \textbf{Spatial and Angular Resolution:} All input LF images are uniformly cropped and resized to a spatial resolution of $256 \times 256$ pixels. The angular resolution is standardized to $9 \times 9$ (81 views) to maintain consistent epipolar geometry across different sources.
2. \textbf{Ground Truth Normalization:} Raw GT disparity values exhibit significant scale variations across different datasets. To mitigate this, we apply a per-scene 2nd-to-98th percentile clipping to the GT disparity maps to eliminate extreme outliers caused by rendering artifacts, followed by min-max normalization to the range $[0, 1]$.
3. \textbf{Domain-Balanced Sampling:} Given the inherent imbalance in the number of scenes across different reflectance domains (Lambertian vs. Non-Lambertian), we employ a `WeightedRandomSampler` during the data loading phase. The sampling weights are inversely proportional to the frequency of each domain, ensuring that the network receives balanced gradient updates from minority domains (e.g., pure specular scenes).
4. \textbf{Data Augmentation:} To improve generalization and prevent overfitting, we apply random horizontal flipping. Importantly, this spatial transformation is consistently applied across all 81 angular views. To preserve the underlying physical parallax constraints and EPI structural integrity, the angular indices along the horizontal dimension are simultaneously inverted (i.e., mapping view $(u, v)$ to $(-u, v)$ when the spatial coordinate $x$ is flipped), ensuring that the slope of the horizontal EPI lines remains physically consistent.

\paragraph{Dataset Summary.}

Table I summarizes the key characteristics of the datasets utilized in our experiments.

\textbf{TABLE I}  
\textbf{SUMMARY OF THE LIGHT FIELD DATASETS USED IN OUR EXPERIMENTS}

\begin{table*}[!t]
\small
\caption{Comparison of performance metrics.}
\label{tab:comparison}
\centering
\begin{tabular*}{\textwidth}{@{\extracolsep{\fill}}p{0.13\textwidth}p{0.13\textwidth}p{0.13\textwidth}p{0.13\textwidth}p{0.13\textwidth}p{0.13\textwidth}}
\toprule
Dataset \& Scene Type \& Angular Res. \& Spatial Res. (Original) \& GT Source \& Train / Val Scenes \\
\midrule
HCInew \& Lambertian / Mixed \& $9 \times 9$ \& $512 \times 512$ \& Synthetic (Ray-tracing) \& 50 / 10 \\
Wanner HCI \& Lambertian \& $9 \times 9$ \& $512 \times 512$ \& Synthetic (Rendering) \& 30 / 5 \\
Non-Lambertian (Zhenglong) \& Non-Lambertian \& $9 \times 9$ \& $512 \times 512$ \& Synthetic (BRDF-based) \& 80 / 15 \\
UrbanLF-Syn \& Mixed (Urban) \& $9 \times 9$ \& $1024 \times 1024$ \& Synthetic (3D City Models) \& 49 / 6 \\
\textbf{Total} \& \textbf{All Domains} \& \textbf{$9 \times 9$} \& \textbf{$256 \times 256$ (Unified)} \& \textbf{-} \& \textbf{209 / 36} \\
\bottomrule
\end{tabular*}
\end{table*}
\textit{Note: The spatial resolution for all datasets is uniformly preprocessed to $256 \times 256$ prior to network input.}

\subsubsection{Implementation Details}

\textbf{Hardware and Software Environment.} 
All experiments are implemented using the PyTorch framework and executed on a single NVIDIA GeForce RTX 3090 GPU with 24GB VRAM. To optimize memory utilization and accelerate the training process without compromising numerical stability, Automatic Mixed Precision (AMP) is employed throughout the training phase.

\textbf{Dataset and Preprocessing.} 
We evaluate our proposed method on four widely-used light field datasets: HCInew [1], Wanner_HCI [2], Non-lambertian_zhenglong [3], and UrbanLF-Syn [4]. The HCI-Old dataset is excluded to ensure format compatibility and consistency. In total, the curated dataset comprises 209 training scenes and 36 validation scenes, encompassing Lambertian (pure diffuse), Non-Lambertian (specular/scattering), and Mixed (urban) scenarios. For uniform processing, all input light field images are resized to a spatial resolution of $256 \times 256$ with an angular resolution of $9 \times 9$ (81 views). The ground-truth disparity maps are truncated at the 2nd and 98th percentiles on a per-scene basis to eliminate extreme outliers, and subsequently normalized to the range of $[0, 1]$.

\textbf{Training Hyperparameters.} 
The network is optimized using the AdamW optimizer with an initial learning rate of $1 \times 10^{-4}$ and a weight decay of $1 \times 10^{-4}$. We employ a Cosine Annealing learning rate scheduler to smoothly decay the learning rate to a minimum of $1 \times 10^{-6}$ over the course of training. Due to the high memory footprint of processing 81-view light fields, the mini-batch size is set to 1, and we utilize gradient accumulation over 4 steps, yielding an effective batch size of 4. The model is trained for 100 epochs. The key hyperparameters are summarized in Table I.

\textbf{Data Augmentation and Sampling Strategy.} 
To mitigate overfitting and improve generalization, we apply random horizontal flipping as a spatial data augmentation. Importantly, this transformation is consistently applied across all 81 angular views. To strictly preserve the epipolar geometry, the horizontal angular indices are simultaneously inverted alongside the spatial horizontal flipping, ensuring the physical consistency of EPI slopes. Furthermore, to address the inherent domain imbalance among Lambertian, Non-Lambertian, and Mixed scenes, we implement a domain-balanced inverse-frequency weighted sampling strategy using `WeightedRandomSampler` during the construction of the training data loader.

\textbf{Loss Function Configuration.} 
The overall objective function is a weighted sum of the primary depth estimation loss and the physical reflection regularization loss. Specifically, the depth estimation loss $\mathcal{L}_{depth}$ is formulated as the Smooth L1 loss between the predicted disparity map $\hat{D}$ and the ground truth $D$:
\begin{equation}
\label{eq:eq35}
\mathcal{L}_{depth} = \frac{1}{N} \sum_{i=1}^{N} \text{SmoothL1}(\hat{D}_i, D_i)
\end{equation}
The physical reflection regularization loss $\mathcal{L}_{ref}$ enforces the consistency of the dual-mask physical model by penalizing deviations from the assumed BRDF priors. The total loss is defined as:
\begin{equation}
\label{eq:eq36}
\mathcal{L}_{total} = \lambda_{depth} \mathcal{L}_{depth} + \lambda_{ref} \mathcal{L}_{ref}
\end{equation}
where the weighting hyperparameters $\lambda_{depth}$ and $\lambda_{ref}$ are empirically set to 1.0 and 0.1, respectively.

\subsubsection{Evaluation Metrics}

Based on the metrics defined in [5], we evaluate the depth estimation performance using Mean Squared Error (MSE) and the percentage of Bad Pixels (BadPix). The MSE measures the average squared difference between the predicted disparity $\hat{D}$ and the ground truth $D$:
\begin{equation}
\label{eq:eq37}
MSE = \frac{1}{N} \sum_{i=1}^{N} (\hat{D}_i - D_i)^2
\end{equation}
The BadPix metric calculates the proportion of pixels where the absolute disparity error exceeds a predefined threshold $\tau$ (set to 0.07 in our experiments):
\begin{equation}
\label{eq:eq38}
BadPix = \frac{1}{N} \sum_{i=1}^{N} \mathbb{I}(|\hat{D}_i - D_i| > \tau)
\end{equation}
where $\mathbb{I}(\cdot)$ is the indicator function. Lower values for both metrics indicate superior depth estimation accuracy.

\subsubsection{Quantitative and Qualitative Comparison}

\textbf{Quantitative Results.} 
We compare our Unified Dual-Mask Physical Model against several state-of-the-art (SOTA) light field depth estimation methods, including EPINET, LFattNet, and DistgSSR. As demonstrated in our comprehensive evaluations, our method achieves the lowest MSE and BadPix scores across all evaluated datasets. Notably, on the Non-Lambertian dataset, our Dual-Mask Physical Model significantly reduces the Non-Lambertian MSE by 18.4\% compared to the second-best baseline, demonstrating its superior capability in handling complex reflectance properties without compromising geometric accuracy.

\textbf{Qualitative Results.} 
Visual comparisons further corroborate the quantitative findings. Traditional EPI-based methods often produce severe depth artifacts and blurring in specular and scattering regions due to the violation of the photo-consistency assumption. In contrast, our method yields sharp and accurate depth boundaries even in highly reflective areas. The dual-mask mechanism successfully isolates the specular highlights from the underlying diffuse geometry, resulting in visually coherent disparity maps that closely match the ground truth, particularly in complex mixed urban scenes.

\subsubsection{Ablation Study}

To verify the effectiveness of the proposed components, we conduct a comprehensive ablation study on the validation set. We evaluate three variants of our model: (1) the baseline network without physical priors, (2) the baseline with a single-mask reflection model, and (3) the full Unified Dual-Mask Physical Model. 

The results indicate that incorporating physical priors is essential for unified depth estimation across diverse reflectance domains. Transitioning from the baseline to the single-mask variant yields a moderate improvement in Lambertian scenes but struggles with mixed urban environments. However, the full dual-mask configuration achieves the best overall performance. This breakthrough is substantially driven by the \textbf{Physical Reflection Regularization}, which explicitly decouples the view-dependent specular layer from the view-independent diffuse layer. Consequently, the network learns a more robust representation of the underlying geometric structure, confirming the necessity and efficacy of the proposed physical model.